Brightness/contrast jitter increased mean marked accuracy by 1.13 percentage points. Its effect on movement flips was small (6.23\% to 6.06\%). In the mark-free controls it also improved mean clean accuracy, from 73.53\% to 74.92\%. Higher digit accuracy can therefore occur without exposure to the synthetic background feature. Gaussian weighting alone yielded 71.42\% marked accuracy and 5.19\% movement flips. The ordinary and jitter models both allowed a linear probe to classify the symbol at about chance.

Translation gave the strongest standalone result. Marked ACC rose by $13.12\pm2.98$ points and movement flips fell by $2.82\pm0.76$ points versus the matched baseline. It also raised ACC by $10.34\pm3.80$ points with no marks in training or testing. Both gains occurred in all five seeds. When those clean-trained models first encountered marks at test time, their movement flips averaged 11.69\% after translation training and 10.50\% without it. The clean accuracy gain alone therefore gives no guarantee of tolerance to an unfamiliar mark.

Adding Gaussian weighting to translation yielded 79.93\% marked ACC and 3.66\% movement flips, compared with 80.63\% and 3.41\% for translation alone. Gaussian weighting helped rotation on both outcomes. Across the five standalone conditions, the linear shape probe ranged from 50.13\% to 50.26\%, and all ten test clusters were occupied in every run.
